% ===========================================================================
%  SCX——
%  SCX and Philosophy of Education — Educational Implications of the Economy Theorem
%  Target:  /  
%  Date: 2026-06-30
% ===========================================================================

\documentclass[11pt,a4paper]{article}

% ---- Chinese & Font Support (XeLaTeX) ----
\usepackage[UTF8,fontset=none]{ctex}
\setCJKmainfont{SimSun}[BoldFont=SimHei,ItalicFont=KaiTi]
\setCJKsansfont{SimHei}
\setCJKmonofont{KaiTi}
\setmainfont{Times New Roman}

% ---- Mathematics ----
\usepackage{amsmath,amssymb,amsthm,mathtools}
\usepackage{bm}
\usepackage{bbm}
\usepackage{dsfont}

% ---- Theorem Environments ----
\newtheorem{theorem}{Theorem}[section]
\newtheorem{proposition}[theorem]{Proposition}
\newtheorem{corollary}[theorem]{Corollary}
\newtheorem{lemma}[theorem]{Lemma}
\newtheorem{definition}[theorem]{Definition}
\newtheorem{assumption}[theorem]{Axiom}
\theoremstyle{remark}
\newtheorem{remark}[theorem]{Remark}
\newtheorem{example}[theorem]{Example}

% ---- Layout ----
\usepackage[margin=2.5cm]{geometry}
\usepackage{booktabs}
\usepackage{graphicx}
\usepackage{hyperref}
\usepackage{enumitem}
\usepackage[capitalise,noabbrev]{cleveref}

% ---- Math Operators ----
\DeclareMathOperator{\PE}{PE}
\DeclareMathOperator{\Var}{Var}
\DeclareMathOperator*{\argmin}{arg\,min}
\DeclareMathOperator*{\argmax}{arg\,max}
\DeclareMathOperator{\KL}{D_{KL}}

% ---- Custom notations ----
\newcommand{\R}{\mathbb{R}}
\newcommand{\N}{\mathbb{N}}
\newcommand{\Y}{\mathcal{Y}}
\newcommand{\X}{\mathcal{X}}
\newcommand{\Sstates}{\mathcal{S}}
\newcommand{\Ppos}{\mathcal{P}}
\newcommand{\E}{\mathbb{E}}
\newcommand{\V}{\mathbb{V}}
% \I removed due to package conflict (dsfont vs CJK)

% ---- Honest critique markers ----
\newcommand{\rigorous}{\textnormal{\textbf{[]}}}
\newcommand{\heuristic}{\textnormal{\textbf{[]}}}
\newcommand{\openproblem}{\textnormal{\textbf{[]}}}
\newcommand{\metaphor}{\textnormal{\textbf{[]}}}

\title{\textbf{\LARGE SCX\\[6pt] ——}}

\author{%
SCX Philosophy Working Group\thanks{SCX
SCX2026}\\
\textit{ /  / }
}

\date{\today}

\begin{document}

\maketitle

% ===========================================================================
% ABSTRACT
% ===========================================================================
\begin{abstract}
SCXState-Conditioned eXpertise
——Robbins-MonroSpring SE-1Chernoff-Hoeffding
Theorem 1$I(Y; P \mid S) > 0$——

SCX\textbf{}``''
——$\Delta_s > 0$
\textbf{}bug
\textbf{}Spring SE-1
$t_{\text{crit}}$
PL$O(1/t)$``''
\textbf{}SCX Theorem 1
——$F_1$
$1 - \exp(-2M\Delta_s^2)$
\textbf{}$\Sstates_i$
$d_i$$\kappa_i$PL$\mu_i$
$\theta_i^*(t)$``''——$\Omega(\max_i d_i)$


``''



\textbf{}SCX
Robbins-MonroSpring
\end{abstract}

% ===========================================================================
% SECTION 1: INTRODUCTION
% ===========================================================================
\section{}

\subsection{SCX}

SCX\cite{scx2026theorems,scx2026cc_audit}State Crystallization
SitusSpringYajieCercis


\begin{enumerate}[label=(\arabic*)]
    \item \textbf{} $\Sstates = \{s_1, \ldots, s_N\}$$N$``''
          $s_i \in \R^{d_s}$/
    \item \textbf{} $\Delta_s(t) = p_{\text{noisy},s}(t) - p_{\text{clean},s}(t)$
          $s$$t$$M$
          $\Delta_s > 0$
    \item \textbf{SpringSE-1}$\theta_t$Robbins-Monro
          \begin{equation}
              \theta_{t+1} = \theta_t - \alpha_t \cdot \nabla\mathcal{L}_{\text{Spring}}(\theta_t; \mathcal{B}_t) + \xi_t
          \end{equation}
          $\alpha_t$$\sum \alpha_t = \infty$$\sum \alpha_t^2 < \infty$
          $\xi_t$
    \item \textbf{Theorem 1}$M$$F_1$
          \begin{equation}
              F_1 \geq 1 - \frac{1}{\eta}\sum_{s \in \Sstates} \rho_s \cdot \exp(-2M\Delta_s^2)
          \end{equation}
    \item \textbf{}$P$$S$$Y$
          $I(Y; P \mid S) > 0$SCX
\end{enumerate}

\subsection{}

——\textbf{}
Spring$\theta_t$\textbf{}$\Sstates$\textbf{}
Yajie$\{E_m\}_{m=1}^M$\textbf{}$\Delta_s(t)$SCX
normative force

——\textbf{}structural isomorphism~\ref{tab:mapping}


\begin{table}[htbp]
\centering
\caption{SCX}
\label{tab:mapping}
\begin{tabular}{@{}lll@{}}
\toprule
\textbf{SCX} & \textbf{} & \textbf{} \\
\midrule
 $s_i$ & $\R^{d_s}$ & / \\
Spring $\theta_t$ & $\R^{|\theta|}$ &  \\
 $\Delta_s(t)$ & $[-1, 1]$ & $s$ \\
SE-1 & $\theta_{t+1} = \theta_t - \alpha_t \nabla\mathcal{L}$ & / \\
 $\mathcal{M}_t$ & $\{(s_i, \text{score}_t(s_i))\}$ &  \\
 $E_m$ & $\X \to \{0, 1\}$ & / \\
Cercis & $[0, 1]$ &  \\
\bottomrule
\end{tabular}
\end{table}

\subsection{}

``---''SCX



% ===========================================================================
% SECTION 2: 
% ===========================================================================
\section{}

\subsection{}

\begin{definition}[——SCX]
\label{def:economy}
Spring$t$$t' > t$$\theta_{t'} = \theta_t$
$\mathcal{M}_{t'} = \mathcal{M}_t$
$s$$\Delta_s$$\tau \to \infty$
\begin{equation}
    \boxed{
    \lim_{\tau \to \infty} \Delta_s(t + \tau) = 0
    }
\end{equation}
\textbf{}
\end{definition}

``''——


\subsection{}

\begin{theorem}[——]
\label{thm:economy}
$\theta_t$Spring$\mathcal{M}_t$``''
$\tau \geq 0$$\theta_{t+\tau} = \theta_t$$\mathcal{M}_{t+\tau} = \mathcal{M}_t$
$\{E_m\}_{m=1}^M$\textbf{}$\gamma > 0$
$m$$s$$q_m(s, \tau)$$\ell_1$

\begin{equation}
    q_m(s, \tau+1) \leq (1 - \gamma) \cdot q_m(s, \tau)
\end{equation}

\begin{equation}
    \boxed{
    \lim_{\tau \to \infty} \Delta_s(t + \tau) = 0
    }
\end{equation}
\end{theorem}

\begin{proof}
\textbf{1}$E_m$$h_i = \phi_{\theta_t}(s_i) + \PE(p_i)$
$\theta_t$$\phi_{\theta_t}$$q_m$——

calibration$q_m(s, \tau+1) \leq (1-\gamma)q_m(s,\tau)$
``''

\textbf{2}$M$$(1-\gamma)^\tau$
\textbf{}$s$
\begin{align}
    \lim_{\tau \to \infty} p_{\text{clean},s}(t+\tau)
    &= \lim_{\tau \to \infty} \E[\mathbf{1}[E_m(s) \neq y] \mid s \text{ }] \nonumber\\
    &= P(\text{}) = \frac{|\Y| - 1}{|\Y|}\\
    \lim_{\tau \to \infty} p_{\text{noisy},s}(t+\tau)
    &= \lim_{\tau \to \infty} \E[\mathbf{1}[E_m(s) \neq y] \mid s \text{ }] \nonumber\\
    &= P(\text{}) = \frac{|\Y| - 1}{|\Y|}
\end{align}

\textbf{3}
\begin{equation}
    \lim_{\tau \to \infty} \Delta_s(t+\tau)
    = \frac{|\Y| - 1}{|\Y|} - \frac{|\Y| - 1}{|\Y|} = 0
\end{equation}

Chernoff-Hoeffding$F_1 \geq 1 - \frac{1}{\eta}$——vacuous bound


\textbf{4$\gamma$}
$\Delta_s(t+\tau) \leq \Delta_s(t) \cdot (1-\gamma)^\tau$
$T_{1/2} = \log 2 / \log(1/(1-\gamma)) \approx (\log 2)/\gamma$$\gamma$
$\gamma = 0.1$10\%$T_{1/2} \approx 6.6$
\end{proof}

\rigorous{} \textbf{}1--4
\textbf{}$q_m(s, \tau+1) \leq (1-\gamma)q_m(s,\tau)$
Ebbinghaus, 1885
Wixted \& Carpenter, 2007
$\gamma$$s$——``''

\subsection{}

\begin{corollary}[——]
\label{cor:spaced_repetition}
$\theta_t$$s$

\textbf{$s$}
$T_{\text{crit}} = O(1/\gamma_s)$$s$$\Delta_s$
``''

\textbf{spaced repetition——
$\Delta_s > 0$}``''
\end{corollary}

\textbf{}

\begin{enumerate}[label=(\arabic*)]
    \item \textbf{``''}——
          
          ``''
          
    \item \textbf{$\gamma_s$}$\gamma_s$$s$
          ——$s$
          $\Delta_s(0)$$s$
          $\gamma_s$
          $\gamma_s$
    \item \textbf{``''}$\tau = 1$
          ——$\Delta_s$
          $\Delta_s$\textbf{}
    \item \textbf{bug}$\gamma_s$
          \textbf{}——$\gamma_s$
          $\gamma_s = 0$——
          ``''
\end{enumerate}

\heuristic{} \textbf{}``''$q_m \to 0$
``''
——
\S6

% ===========================================================================
% SECTION 3: 
% ===========================================================================
\section{Spring}

\subsection{Spring SE-1}

Spring SE-1(Spring Self-Evolution 1Robbins-Monro)
\textbf{}

\begin{enumerate}[label=\,\arabic*:, leftmargin=*]
    \item \textbf{$t < t_{\text{crit}}$}$\theta_t$$K!$
          $\xi_t$
          $\|\nabla\mathcal{L}(\theta_t)\|$——
          SCX\cite{scx2026spring_convergence}1.2
          $\min_{\tau \leq T} \E[\|\nabla\mathcal{L}(\theta_\tau)\|^2] = O(\log T / \sqrt{T})$
    \item \textbf{$t \geq t_{\text{crit}}$}``''
          Polyak-ŁojasiewiczPL$O(1/t)$——
          1.1$\E[\mathcal{L}(\theta_t) - \mathcal{L}(\theta^*)] \leq \frac{2L\sigma^2}{\mu^2} \cdot \frac{1}{t}$
\end{enumerate}

\begin{definition}[——]
\label{def:critical_point}
$t_{\text{crit}}$Spring SE-1
\begin{equation}
    \boxed{
    t_{\text{crit}} = \inf\left\{t : \E[\|\nabla\mathcal{L}(\theta_t)\|^2] \leq \epsilon_{\text{basin}}
    \text{  } \lambda_{\min}(\nabla^2\mathcal{L}(\theta_t)) \geq -\sqrt{\rho\epsilon_{\text{basin}}}\right\}
    }
\end{equation}
$\epsilon_{\text{basin}} > 0$``''``''
$\rho > 0$
\end{definition}

\subsection{}

\begin{theorem}[]
\label{thm:critical_point}
Spring$\mathcal{L}$$L$-$[\theta^*]$
$\Theta/S_K$PLPL$\mu > 0$
$t_{\text{crit}} < \infty$$t \geq t_{\text{crit}}$
\begin{equation}
    \boxed{
    \E[\mathcal{L}(\theta_t) - \mathcal{L}(\theta^*)] = O\left(\frac{1}{t}\right)
    }
\end{equation}
$t < t_{\text{crit}}$$\min_{\tau \leq t} \E[\|\nabla\mathcal{L}(\theta_\tau)\|^2] = O(\log t / \sqrt{t})$
\end{theorem}

\begin{proof}[]
\textbf{1——}SCX1.2
\begin{equation}
    \min_{1 \leq \tau \leq t} \E[\|\nabla\mathcal{L}(\theta_\tau)\|^2]
    \leq \frac{\mathcal{L}(\theta_0) - \mathcal{L}_{\inf} + \frac{L\alpha^2\sigma^2}{2}H_t}{\alpha \cdot (2\sqrt{t} - 1)}
    = O\left(\frac{\log t}{\sqrt{t}}\right)
\end{equation}
``''``/''

\textbf{2}SCX1.3Hessian
$\E[T_{\text{escape}}] = O(\log d_{\text{eff}} / (\rho \alpha_t \sigma^2))$


\textbf{3——}$\theta_t$$\theta^*$
$\mathcal{N}$PL$\mathcal{N}$$\alpha_t = \frac{2}{\mu(t + t_0)}$
1.1$O(1/t)$

\textbf{4}1PL
1$t_{\text{crit}}$$\theta_t$
$t_{\text{crit}} = O(K! \cdot \log d_{\text{eff}} / (\rho \alpha \sigma^2))$
$K!$\hfill $\square$
\end{proof}

\rigorous{} \textbf{1--3SCX1.11.21.3
4PLSpring1
$t_{\text{crit}}$$K!$
}

\subsection{}

\begin{corollary}[——]
\label{cor:qualitative_change}
Spring SE-1\textbf{}$t_{\text{crit}}$

\begin{enumerate}[label=(\arabic*)]
    \item \textbf{$t < t_{\text{crit}}$/}——
          $\Delta_s(t)$SCX3.1
          ``''——
          
    \item \textbf{$t \geq t_{\text{crit}}$/}
          ``''$O(1/t)$
          ——$\Delta_s(t)$PL
          ``''——\textbf{}\textbf{}
\end{enumerate}
\end{corollary}

\textbf{}

\begin{enumerate}[label=(\arabic*)]
    \item \textbf{``''}$t < t_{\text{crit}}$
          \textbf{}——
          ``''
          \textbf{}
          \textbf{}——
          
    \item \textbf{``''}``''Spring
          $\theta_t$PLPL$O(1/t)$
          $O(\log t/\sqrt{t})$\textbf{}——
          ``''
          \textbf{}\textbf{}——
          Robbins-Monro$\alpha_t \to 0$
    \item \textbf{}$t_{\text{crit}}$$O(K!)$——
          $K$``''
          $K=4$$4! = 24$——
          1$1! = 1$——
          
    \item \textbf{``''}$\Delta_s(t)$
          SCX3.1→→
          ``''——
          
\end{enumerate}

\heuristic{} \textbf{}$t_{\text{crit}}$$K!$——
``''$K$PL
PL
Spring$\theta$——

Spring SE-1

% ===========================================================================
% SECTION 4: 
% ===========================================================================
\section{}

\subsection{SCX Theorem 1F1}

SCXTheorem 1$F_1$

\begin{theorem}[SCX Theorem 1——F1]
\label{thm:scx1}
$M$$\{E_1, \ldots, E_M\}$$s$$v_m(s) = \mathbf{1}[E_m(s) \neq y]$
$H_0$$v_m \sim \text{Bernoulli}(p_{\text{clean},s})$$H_1$
$v_m \sim \text{Bernoulli}(p_{\text{noisy},s})$$p_{\text{noisy},s} > p_{\text{clean},s}$
$\Delta_s = p_{\text{noisy},s} - p_{\text{clean},s} > 0$
Chernoff-Hoeffding
\begin{equation}
    \boxed{
    F_1 \geq 1 - \frac{1}{\eta} \sum_{s \in \Sstates} \rho_s \cdot
    \exp\!\left(-2M \Delta_s^2\right)
    }
\end{equation}
$\eta = \sum_{s} \rho_s \cdot \mathbf{1}[\text{ } s \text{ }]$
$\rho_s$$s$
\end{theorem}

\textbf{}$M$$\Delta_s$$F_1$\textbf{}
$1 - \exp(-2M\Delta_s^2)$1$M=1$$F_1$
$1 - \frac{1}{\eta}\sum_s \rho_s \exp(-2\Delta_s^2)$——

\subsection{}

\begin{corollary}[]
\label{cor:multi_expert}
SCX Theorem 1

\begin{enumerate}[label=(\arabic*)]
    \item \textbf{}$M=1$
          Theorem 1$F_1$$1 - \frac{1}{\eta}\sum_s \rho_s \exp(-2\Delta_s^2)$——
          $\Delta_s$$\exp(-2\Delta_s^2) \approx 1$
          $F_1$1\textbf{``''
          ``''}
    \item \textbf{$F_1$}$M$
          $F_1$$1 - \exp(-2M\bar{\Delta}^2)$1
          $\bar{\Delta} = \min_s \Delta_s$$M \geq 3$$\bar{\Delta} \geq 0.3$
          $F_1 \geq 0.95$——\textbf{``''}
    \item \textbf{``''}Theorem 1
          $v_m$
          ``''$M_{\text{eff}} < M$——Chernoff$M$
          $M_{\text{eff}}$$M$
          $M_{\text{eff}} \approx 1$——
    \item \textbf{SATGRE}
          $\Sstates$
          $\Delta_s$——\textbf{}
          
          $M_{\text{eff}} = 1$Theorem 1$F_1$
          vacuous
\end{enumerate}
\end{corollary}

\textbf{}

\begin{enumerate}[label=(\arabic*)]
    \item \textbf{}$M=1$
          Theorem 1\textbf{}
          ——$F_1$$\exp(-2M\Delta_s^2)$
          formative assessment
    \item \textbf{peer assessment}$M$
          $F_1$$M$
          ——rubric
          $M_{\text{eff}}$
    \item \textbf{——}$M=4$
          
          
          Theorem 1$F_1$
          $\Delta_s$——+
          ``''Dunning-Kruger
    \item \textbf{}Theorem 1$F_1$
          \textbf{}
          \begin{itemize}
              \item \textbf{}→ 
              \item \textbf{}$M_{\text{eff}}$→ 
              \item \textbf{}→ 
          \end{itemize}
          ``+''
          Theorem 1$F_1$
\end{enumerate}

\rigorous{} \textbf{Theorem 1Chernoff-Hoeffding
(1) ——
(2) $\Delta_s$——
(3) 
Theorem 1}

% ===========================================================================
% SECTION 5: 
% ===========================================================================
\section{——}

\subsection{}

SCX\textbf{}$\Sstates$
$i$\textbf{}$\Sstates_i$

\begin{definition}[]
\label{def:individual_space}
$i$$\Sstates_i \subset \R^{d_i}$
\begin{enumerate}[label=(\arabic*)]
    \item \textbf{$d_i$}$i$$d_i$
          ——
    \item \textbf{$\kappa_i$}``''——$\kappa_i$
          
    \item \textbf{PL$\mu_i$}$\mu_i$
          ``''
    \item \textbf{$\boldsymbol{\gamma}_i$}$s$$i$
          ——$s$$\Sstates_i$
\end{enumerate}
\end{definition}

\subsection{}

\begin{theorem}[——]
\label{thm:individualized}
$N$$\{1, \ldots, N\}$$\Sstates_i$$d_i$PL$\mu_i$
``''$\pi_{\text{uniform}}$
$\alpha_t = \alpha / \sqrt{t}$
$i$$\theta_i^*$

\begin{equation}
    \boxed{
    \max_{i \in [N]} \mathbb{E}[R_T^{(i)}(\pi_{\text{uniform}})]
    \geq \Omega\!\left(\max_i d_i \cdot \frac{\log T}{\sqrt{T}} + \frac{1}{\min_i \mu_i} \cdot \frac{1}{T}\right)
    }
\end{equation}
\textbf{}$\Sstates_i$\textbf{}
$\mu_i$PL
\end{theorem}

\begin{proof}
\textbf{1}SCX1.2
$\min_{\tau \leq T} \E[\|\nabla\mathcal{L}(\theta_\tau)\|^2] = O(\log T / \sqrt{T})$
$d_i$——
$d_i$
$\E[\|\nabla\mathcal{L}(\theta_\tau)\|^2] \geq \Omega(d_i \cdot \log T / \sqrt{T})$
$d_i$

\textbf{2PL}PL$O(1/\mu_i t)$$\mu_i$
``''——$\alpha_t$
$\alpha_t^{(i)} = 2/(\mu_i(t+t_0^{(i)}))$——$\mu_i$$t_0^{(i)}$

\textbf{3}$\propto d_i \log T / \sqrt{T}$
$\propto 1/(\mu_i T)$
$\max_i d_i / \bar{d}$$\max_i (1/\mu_i) / \overline{(1/\mu)}$
PL\hfill $\square$
\end{proof}

\rigorous{} \textbf{SCX1.1--1.2
}1.4$K!$
$K_i$——$K_i$

\subsection{}

\begin{corollary}[——]
\label{cor:individualized}
Spring SE-1

\begin{enumerate}[label=(\arabic*)]
    \item \textbf{}$d_i$$\mu_i$
          $\Omega(\max_i d_i)$
          ``''
          $\propto d_{\max}$
    \item \textbf{``''}SCX``''
          $\mu_i$——
          ``''``''——
          $\alpha_t^{(i)} = 2/(\mu_i t)$
    \item \textbf{}$\alpha_t^{(i)}$PL$\mu_i$
          $\mu_i$
          $\mu_i$+
          ``''$\alpha_t$
          ``''$\alpha_t$$\mu_i$→
          ``''$\alpha_t$$\mu_i$→
    \item \textbf{}$i$$d_i$——
          ``''
          ``''——
          
\end{enumerate}
\end{corollary}

\textbf{}

\begin{enumerate}[label=(\arabic*)]
    \item \textbf{}$\Sstates_i$
          ——$\mu_i$$\gamma_{s,i}$
          ``''
    \item \textbf{}Khan AcademyALEKS
          $\Delta_s(t)$——
          $\Sstates_i$$\alpha_t^{(i)}$
          Theorem \ref{thm:individualized}
    \item \textbf{}``''————
          SCX$\Sstates_i$
          $\{\theta_t^{(i)}\}_{t=0}^{T}$
          ``''——
          ``$\Omega(\max_i d_i)$''
\end{enumerate}

\heuristic{} \textbf{}$\Sstates_i$$d_i$, $\mu_i$, $\kappa_i$, $\gamma_{s,i}$
PL$\mu_i$——
``''
Spring SE-1

% ===========================================================================
% SECTION 6: 
% ===========================================================================
\section{}

SCX
\rigorous{}\heuristic{}
\metaphor{}

\subsection{→}

\begin{table}[htbp]
\centering
\caption{}
\label{tab:economy_audit}
\begin{tabular}{@{}p{3.5cm}p{6cm}p{5cm}@{}}
\toprule
\textbf{} & \textbf{SCX} & \textbf{} \\
\midrule
 &
$q_m(s, \tau+1) \leq (1-\gamma)q_m(s,\tau)$ &
+``''\\
\midrule
$\gamma$ &
$\gamma$$s \in \Sstates$ &
 vs  vs \\
\midrule
 &
$q_m \to 0 \implies \Delta_s \to 0$ &
\\
\midrule
/ &
 &
$\Delta_s$\\
\bottomrule
\end{tabular}
\end{table}

\textbf{}``''\textbf{}——
``''
\textbf{}savings effect——

SCX\textbf{}$r_s(\tau)$——$\Delta_s \to 0$
$r_s > 0$``''——
``''

\subsection{→PL}

\textbf{PL}Polyak-Łojasiewicz
$\frac{1}{2}\|\nabla\mathcal{L}(\theta)\|^2 \geq \mu(\mathcal{L}(\theta) - \mathcal{L}(\theta^*))$

Liu et al., 2022\textbf{``''}——
$\mathcal{L}_{\text{human}}(\theta)$``
''Spring SE-1
\textbf{}

$t_{\text{crit}} = O(K! \cdot \log d_{\text{eff}} / (\rho \alpha \sigma^2))$
$K=5$
$120 \times$——


\subsection{→}

\textbf{}

\begin{enumerate}[label=(\arabic*)]
    \item \textbf{}
          ``''
          $M_{\text{eff}}$$M_{\text{nominal}}$
    \item \textbf{}``''``''``''
          ——
          Theorem 1Chernoff\textbf{}——
          $F_1$``$M_{\text{eff}}$''
          ``''
    \item \textbf{}——
          $\Delta_s$testing effect
          \textbf{}$\Delta_s$Theorem 1
          
    \item \textbf{}Theorem 1``''$y$
          
          ``''``''——
          ``''``''
          \textbf{}
\end{enumerate}

\textbf{}``''
``\textbf{}reflective equilibrium''——

$\lim_{M \to \infty} \frac{1}{M}\sum_{m=1}^M v_m(s)$
``''$p_s$

\subsection{→}

\textbf{$\Sstates_i$}$d_i$$\mu_i$PL
$\kappa_i$``''——

\textbf{Theorem 3}——
$\mu_i$$\Delta_s$
/

\subsection{}

\begin{table}[htbp]
\centering
\caption{}
\label{tab:honesty}
\begin{tabular}{@{}llll@{}}
\toprule
\textbf{} & \textbf{SCX} & \textbf{} & \textbf{} \\
\midrule
 &  & \heuristic  &  \\
 & Spring SE-1 & \metaphor  &  \\
$>$ & Theorem 1 & \heuristic  &  \\
 &  & \heuristic  & $\Sstates_i$ \\
\bottomrule
\end{tabular}
\end{table}

\textbf{}SCX\textbf{}
possibility proof——
\textbf{}——
\textbf{}falsifiable hypothesis


% ===========================================================================
% SECTION 7: DISCUSSION
% ===========================================================================
\section{}

\subsection{SCX}

\textbf{SCX}

\begin{enumerate}[label=(\arabic*)]
    \item Spring SE-1——
          
    \item $F_1$$\exp(-2M\Delta_s^2)$Theorem 1——
          Chernoff
    \item →——
          
    \item ——
          
\end{enumerate}

\textbf{SCX}

\begin{enumerate}[label=(\arabic*)]
    \item 
    \item $K$PL$\mu$
    \item ``''Theorem 1
    \item effect size——
          
\end{enumerate}

\subsection{}

\textbf{VygotskyZPD}SCXZPD
$\|\nabla\mathcal{L}(\theta)\|$——
\textbf{}scaffolding
PL——Vygotsky``''


\textbf{BloomMastery Learning}Bloom
80\%SCX
$\Delta_s(t) \geq \tau_{\text{mastery}}$$\Sstates$
$\Delta_s$$\tau_{\text{mastery}}$
$\Delta_s$——

\textbf{SwellerCLT}intrinsic
extraneousgermaneSCX
``''——
Spring SE-1

\subsection{Open Problem}

\begin{enumerate}[label=(\arabic*), leftmargin=*]

    \item \textbf{PL}Duolingo
          Khan AcademyPL$O(1/t)$
          PL

    \item \textbf{$\gamma_s$}
          SCX$s$
          $\gamma_s$

    \item \textbf{$M_{\text{eff}}$}
          
          inter-rater reliability

    \item \textbf{Spring SE-1}Spring``''——
          
          
          SCX``''

    \item \textbf{SCX}——
          SCXSpring
          
          $M$

    \item \textbf{Theorem 3}SCXTheorem 3
          ``''
          ``''——
          ——

\end{enumerate}

\subsection{——}

————
SCX\textbf{}\textbf{}

\textbf{}Ebbinghaus
BloomSCX——


\textbf{}——
$\gamma > 0$PL
$M_{\text{eff}} \geq 3$$\Sstates_i$
——\textbf{}
——

\textbf{}\textbf{}
\textbf{}
SCX\textbf{}Spring SE-1\textbf{}
``''——Spring SE-1
——

% ===========================================================================
% ACKNOWLEDGMENTS
% ===========================================================================
\section*{}

——Spring SE-1Theorem 1F1——
SCX\cite{scx2026theorems,scx2026cc_audit,scx2026spring_convergence,scx2026spring_hostile}
SCX
SCX\cite{scx2026spring_hostile}——


% ===========================================================================
% BIBLIOGRAPHY
% ===========================================================================
\begin{thebibliography}{99}

\bibitem{scx2026theorems}
SCX Project.
\newblock {Theorem 1--4}: Multi-expert consistency, weak feature limits,
  unidentifiability, and minimax optimality.
\newblock Technical report, \texttt{theory/theorems/}, 2026.

\bibitem{scx2026cc_audit}
SCX Project.
\newblock Multi-head spring and positional encoding analysis: {CC} audit report.
\newblock Technical report,
  \texttt{theory/self\_evolution/multi\_head\_spring\_and\_positional\_encoding\_analysis.md},
  2026.

\bibitem{scx2026spring_convergence}
SCX Project.
\newblock Spring .
\newblock Technical report,
  \texttt{theory/self\_evolution/spring\_convergence\_analysis.md}, 2026.

\bibitem{scx2026spring_hostile}
SCX Project.
\newblock Spring .
\newblock Technical report,
  \texttt{theory/self\_evolution/spring\_hostile\_review.md}, 2026.

\bibitem{scx2026ppe_derivation}
SCX Project.
\newblock Physical positional encoding ({PPE}) rigorous derivation in the {SCX}
  framework.
\newblock Technical report,
  \texttt{theory/self\_evolution/ppe\_rigorous\_derivation.md}, 2026.

\bibitem{ebbinghaus1885}
H.~Ebbinghaus.
\newblock \emph{\"{U}ber das Ged\"{a}chtnis: Untersuchungen zur experimentellen
  Psychologie}.
\newblock Duncker \& Humblot, Leipzig, 1885.

\bibitem{wixted2007}
J.~T.~Wixted and S.~K.~Carpenter.
\newblock The Wickelgren power law and the Ebbinghaus savings function.
\newblock \emph{Psychological Science}, 18(2):133--134, 2007.

\bibitem{robbins1951}
H.~Robbins and S.~Monro.
\newblock A stochastic approximation method.
\newblock \emph{The Annals of Mathematical Statistics}, 22(3):400--407, 1951.

\bibitem{auer2002}
P.~Auer, N.~Cesa-Bianchi, Y.~Freund, and R.~E.~Schapire.
\newblock The nonstochastic multiarmed bandit problem.
\newblock \emph{SIAM Journal on Computing}, 32(1):48--77, 2002.

\bibitem{jin2017}
C.~Jin, R.~Ge, P.~Netrapalli, S.~M.~Kakade, and M.~I.~Jordan.
\newblock How to escape saddle points efficiently.
\newblock In \emph{ICML}, 2017.

\bibitem{liu2022}
C.~Liu, L.~Zhu, and M.~Belkin.
\newblock Loss landscapes and optimization in over-parameterized non-linear
  systems and neural networks.
\newblock \emph{Applied and Computational Harmonic Analysis}, 59:85--116, 2022.

\bibitem{vygotsky1978}
L.~S.~Vygotsky.
\newblock \emph{Mind in Society: The Development of Higher Psychological Processes}.
\newblock Harvard University Press, 1978.

\bibitem{bloom1968}
B.~S.~Bloom.
\newblock Learning for mastery.
\newblock \emph{Evaluation Comment}, 1(2):1--12, 1968.

\bibitem{sweller1988}
J.~Sweller.
\newblock Cognitive load during problem solving: Effects on learning.
\newblock \emph{Cognitive Science}, 12(2):257--285, 1988.

\bibitem{kruger1999}
J.~Kruger and D.~Dunning.
\newblock Unskilled and unaware of it: How difficulties in recognizing one's own
  incompetence lead to inflated self-assessments.
\newblock \emph{Journal of Personality and Social Psychology}, 77(6):1121--1134, 1999.

\bibitem{ghadimi2013}
S.~Ghadimi and G.~Lan.
\newblock Stochastic first- and zeroth-order methods for nonconvex stochastic
  programming.
\newblock \emph{SIAM Journal on Optimization}, 23(4):2341--2368, 2013.

\end{thebibliography}

\end{document}
